% =====================================================================
%  5장 논의와 한계 (한국어)
%
%  '본 연구가 주장하지 않는 것'은 4장(결과)이 아니라 이 절에 둔다.
%  분량 설계: 약 0.6쪽.
% =====================================================================

\section{논의와 한계}
\label{sec:discussion}

SPR-IRL이 복원한 것은 세 유형의 선호를 하나의 공통 척도 위에 나란히 놓은 그림이다.
모멘텀 축에서 외국인과 개인이 정반대로 갈리고 같은 대비가 컨텍스트 반응에서
재현되며, 기관은 군집을 제외하면 식별되는 선호를 보이지 않는다. 이 세 결과는 세
유형에 동일한 특징과 동일한 모형 형태를 부여한 뒤 파라미터만 분리 학습해 얻은
것이므로 설계가 심어둔 대비가 아니다.

\subsection{본 연구가 주장하지 않는 것}
\label{sec:notclaim}

프로토콜의 마지막 단계로, 결과가 지지하지 않는 주장을 명시한다.

\textbf{(1) 수익률 예측가능성.} 본 연구의 지표는 수급의 방향과 강도에 대한 재현
오차이며 수익률, 샤프지수, 거래비용을 포함하지 않는다. 표~\ref{tab:performance}에서
어떤 매매 규칙도 도출되지 않는다.

\textbf{(2) 수급 예측 성능의 우위.} 지속성 기준선이 방향 정확도와 상관계수에서
SPR-IRL을 앞선다(\S\ref{sec:oos}). 선호 복원과 예측은 다른 특징 집합을 요구한다.
예비 점검에서 지연 자기행동을 네 번째 특징으로 추가하면 모형이 기준선을
넘어서지만(방향 정확도 $0.646/0.569/0.629$) 개인의 $\alpha^{KOSPI}$ 부호가 반전되고
손실구간 계수가 소멸한다. 즉 \textbf{본 논문이 보고한 거울상 중 일부는 수급
자기상관과 분리되지 않는다.}

\textbf{(3) 인과효과.} 복원된 가중치는 표준화된 관측 특징 사이의 조건부 연관이다.
특징--컨텍스트 상관이 $-0.29$에서 $0.22$ 범위이므로 개별 계수의 크기는 다른 항의
통제 조건에 의존한다.

\textbf{(4) 유형 간 음의 연관의 행동적 해석.} 군집 계수와 거울상 컨텍스트 반응은
시장 청산 제약과 분리되지 않는다(\S\ref{sec:beta}).

\textbf{(5) 처분효과의 검정.} 개인의 손실구간 부호는 문헌과 정합하나 월 블록
부트스트랩을 통과하지 못했고, 평균단가 대용치는 계좌별 실현손익이 아니라 집계
거래로 만든 근사치다.

\textbf{(6) 종목 간 일반화.} 단일 종목 설정은 유형 간 이질성 분리를 위한 통제이며
다른 종목과 시장으로의 확장은 검증되지 않았다.

\subsection{향후 과제}

가장 시급한 것은 지연 자기행동을 특징으로 편입해 자기상관을 통제한 뒤 거울상이
살아남는지 확인하는 일이다. 다음으로 ``기관''을 연기금ㆍ투신ㆍ보험 등 세부 주체로
분해하면 상쇄로 가려진 선호가 드러날 수 있다. 마지막으로 본 논문의 근시안 가정을
푸는 것이 남는다. 명제~\ref{prop:lasso}가 보이듯 한 걸음만 내다보는 설정에서 보상은
정책과 겹치지만, 상태 동역학과 다기간 가치를 도입하면 두 대상이 갈라지고 회귀로는
도달할 수 없는 구조적 복원이 된다. SPR-IRL의 대칭 추정과 검증 프로토콜은 그
확장에서도 그대로 쓰일 수 있다.
